%% Drafted from the mapped corpus by scripts/draft_topics.py.
% Model: gpt-5.6-terra; source characters: 16734.
\section{Дискретни случайни величини и разпределения}

\begin{defn}[Случайна величина]
Върху вероятностно пространство $(\Omega,\mathcal A,\mathbb P)$ случайна величина е измерима функция $X:\Omega\to\mathbb R$, тоест $\{\omega:X(\omega)\le x\}\in\mathcal A$ за всяко $x\in\mathbb R$. Ако $\omega\in\Omega$ е елементарен изход, то стойността $X(\omega)$ зависи от реализирания изход $\omega$.
\end{defn}

\begin{defn}
Случайната величина $X$ се нарича \emph{дискретна}, ако множеството от стойностите, които тя приема, е крайно или изброимо.
\end{defn}

Нека $H_1,H_2,\ldots$ е разлагане на $\Omega$ и нека $x_1,x_2,\ldots$ са различни реални числа. Дискретната случайна величина може да се представи като
\[
X(\omega)=\sum_j x_j I_{H_j}(\omega),
\]
където
\[
I_{H_j}(\omega)=
\begin{cases}
1, & \omega\in H_j,\\
0, & \omega\notin H_j
\end{cases}
\]
е индикаторът на множеството $H_j$.

\begin{defn}
Нека $X$ е дискретна случайна величина, която приема стойностите
\[
x_1,x_2,\ldots .
\]
\emph{Дискретно разпределение на $X$} наричаме съответствието
\[
\begin{array}{c|ccccc}
X & x_1 & x_2 & \cdots & x_j & \cdots\\
\hline
\mathbb{P} & p_1 & p_2 & \cdots & p_j & \cdots
\end{array},
\]
където
\[
p_j=\mathbb{P}(X=x_j).
\]
\end{defn}

Вероятностите в таблицата трябва да удовлетворяват две основни условия:
\[
p_j\geq 0\quad\text{за всяко }j,
\qquad
\sum_j p_j=1.
\]
Първото условие изразява неотрицателността на вероятността, а второто --- нормираността: случайната величина трябва да приеме някоя от възможните си стойности с вероятност $1$.

По-общо, вероятността е функция, дефинирана върху $\sigma$-алгебра $\mathcal{A}$ от събития. За всяко $A\in\mathcal{A}$ е изпълнено
\[
\mathbb{P}(A)\geq 0,
\qquad
\mathbb{P}(\Omega)=1.
\]
Ако $A$ и $B$ са несъвместими събития, то
\[
\mathbb{P}(A\cup B)=\mathbb{P}(A)+\mathbb{P}(B).
\]
От свойствата на вероятността следва и монотонността:
\[
A\subseteq B\quad\Longrightarrow\quad \mathbb{P}(A)\leq \mathbb{P}(B).
\]

\begin{defn}[Математическо очакване]
За дискретна случайна величина с вероятности $p_j=\mathbb{P}(X=x_j)$ математическото очакване е
\[
\mathbb{E}X=\sum_j x_jp_j,
\]
когато редът $\sum_j |x_j|p_j$ е сходящ.
\end{defn} \begin{defn}[Дисперсия]
Дисперсията се определя чрез
\[
\mathbb{D}X=\mathbb{E}\bigl[(X-\mathbb{E}X)^2\bigr]
=\sum_j (x_j-\mathbb{E}X)^2p_j.
\]
При $\mathbb E X^2<\infty$ разкриването на квадрата и линейността на очакването дават равенството
\[
\mathbb{D}X=\mathbb{E}X^2-(\mathbb{E}X)^2.
\]
\end{defn}

\section{Пораждаща функция на дискретна случайна величина}

\begin{defn}
Нека $X$ е неотрицателна целочислена дискретна случайна величина, за която
\[
\mathbb{P}(X=k)=p_k,\qquad k=0,1,2,\ldots .
\]
\emph{Пораждаща функция} на $X$ се нарича функцията
\[
G_X(s)=\mathbb{E}s^X=\sum_{k=0}^{\infty}p_ks^k,
\]
за онези стойности на $s$, за които редът е сходящ.
\end{defn}

Коефициентът пред $s^k$ в реда за $G_X(s)$ е именно вероятността
\[
\mathbb{P}(X=k).
\]
Следователно пораждащата функция кодира цялото дискретно разпределение.

\begin{prop}
Нека $X$ е неотрицателна целочислена случайна величина. За пораждащата й функция са в сила следните свойства, като производните при $1$ са леви граници и за формулата за дисперсията се изисква $\mathbb E X^2<\infty$:
\[
G_X(1)=1,
\]
\[
G_X'(1)=\mathbb{E}X,
\]
а при съществуваща втора производна
\[
\mathbb{D}X=G_X''(1)+G_X'(1)-\bigl(G_X'(1)\bigr)^2.
\]
\end{prop}

\begin{proof}
Действително, от нормираността на разпределението непосредствено се получава
\[
G_X(1)=\sum_{k=0}^{\infty}p_k=1.
\]
Първата производна има вида
\[
G_X'(s)=\sum_{k=1}^{\infty}kp_ks^{k-1},
\]
За $0\le s<1$ почленното диференциране е допустимо вътре в радиуса на сходимост. При $s\uparrow1$ монотонната сходимост на неотрицателните членове дава граничната производна, откъдето следва формулата за математическото очакване. Аналогично,
\[
G_X''(1)=\sum_{k=0}^{\infty}k(k-1)p_k.
\]
Тъй като
\[
k^2=k(k-1)+k,
\]
получаваме
\[
\mathbb{E}X^2=G_X''(1)+G_X'(1),
\]
а след това и формулата за дисперсията.
\end{proof}

Пораждащите функции са удобни, защото чрез едно и също пресмятане дават както вероятностите, така и основните моменти на разпределението.

\section{Биномно разпределение}

\subsection{Схема на Бернули}

Разглеждаме опит с два възможни изхода: успех и неуспех. Нека вероятността за успех е $p$, а вероятността за неуспех е
\[
q=1-p.
\]
Случайна величина $Y$, която приема стойност $1$ при успех и стойност $0$ при неуспех, има разпределение на Бернули. За нея
\[
\mathbb{P}(Y=1)=p,
\qquad
\mathbb{P}(Y=0)=q.
\]
Следователно
\[
\mathbb{E}Y=p,
\qquad
\mathbb{D}Y=pq.
\]

Биномното разпределение възниква при фиксиран брой независими опити на Бернули, когато вероятността за успех е една и съща във всеки опит.

\begin{defn}
Нека са извършени $n$ независими опита на Бернули с вероятност за успех $p$. Ако $X$ е броят на успехите в тези $n$ опита, казваме, че $X$ има \emph{биномно разпределение} с параметри $n$ и $p$ и пишем
\[
X\sim\mathrm{Bin}(n,p).
\]
\end{defn}

Случайната величина $X$ приема стойностите
\[
0,1,\ldots,n.
\]
За да се получат точно $k$ успеха, трябва да се изберат $k$ от $n$-те позиции, в които успехът да настъпи. За всяка такава подредба вероятността е
\[
p^kq^{n-k}.
\]
Броят на възможните избори е $\binom{n}{k}$, затова
\[
\mathbb{P}(X=k)=\binom{n}{k}p^kq^{n-k},
\qquad k=0,1,\ldots,n.
\]

\begin{prop}
Вероятностите на биномното разпределение са нормировани:
\[
\sum_{k=0}^{n}\binom{n}{k}p^kq^{n-k}=1.
\]
\end{prop}

\begin{proof}
По биномната формула на Нютон
\[
\sum_{k=0}^{n}\binom{n}{k}p^kq^{n-k}=(p+q)^n.
\]
Тъй като $p+q=1$, получаваме
\[
(p+q)^n=1.
\]
\end{proof}

\begin{example}
Нека една монета има вероятност $p$ да се падне ези и се хвърля $n$ пъти независимо. Ако $X$ е броят на падналите се ези, то
\[
X\sim\mathrm{Bin}(n,p).
\]
При правилна монета $p=\frac12$.
\end{example}

\begin{example}
Правилен зар се хвърля $5$ пъти. Нека $X$ е броят на получените шестици. Всеки опит има успех с вероятност
\[
p=\frac16,
\]
следователно
\[
X\sim\mathrm{Bin}\left(5,\frac16\right).
\]
Например вероятността да се получат точно две шестици е
\[
\mathbb{P}(X=2)
=
\binom52\left(\frac16\right)^2
\left(\frac56\right)^3.
\]
\end{example}

\subsection{Пораждаща функция, математическо очакване и дисперсия}

За $X\sim\mathrm{Bin}(n,p)$ пораждащата функция е
\[
G_X(s)
=
\sum_{k=0}^{n}\binom{n}{k}p^kq^{n-k}s^k.
\]
Като обединим $p^k$ и $s^k$, получаваме
\[
G_X(s)
=
\sum_{k=0}^{n}\binom{n}{k}(ps)^kq^{n-k}.
\]
Отново по биномната формула на Нютон:
\[
G_X(s)=(q+ps)^n.
\]

Първата производна е
\[
G_X'(s)=np(q+ps)^{n-1}.
\]
Следователно
\[
\mathbb{E}X=G_X'(1)
=np(p+q)^{n-1}=np.
\]

Втората производна е
\[
G_X''(s)=n(n-1)p^2(q+ps)^{n-2}.
\]
При $s=1$:
\[
G_X''(1)=n(n-1)p^2.
\]
Затова
\[
\begin{aligned}
\mathbb{D}X
&=G_X''(1)+G_X'(1)-\bigl(G_X'(1)\bigr)^2\\
&=n(n-1)p^2+np-n^2p^2\\
&=np(1-p)=npq.
\end{aligned}
\]

\begin{keythm}
Ако
\[
X\sim\mathrm{Bin}(n,p),
\]
то
\[
\mathbb{E}X=np,
\qquad
\mathbb{D}X=npq=np(1-p).
\]
\end{keythm}

\begin{proof}
Същите формули следват и от представянето
\[
X=X_1+\cdots+X_n,
\]
където $X_i$ са независими случайни величини на Бернули с параметър $p$. Тогава
\[
\mathbb{E}X
=
\mathbb{E}X_1+\cdots+\mathbb{E}X_n
=np,
\]
а по независимостта
\[
\mathbb{D}X
=
\mathbb{D}X_1+\cdots+\mathbb{D}X_n
=npq.
\]
\end{proof}

\section{Геометрично разпределение}

\subsection{Дефиниция и интерпретация}

Геометричното разпределение описва повтаряне на независими опити на Бернули до настъпването на първия успех. Броят на опитите не е предварително ограничен.

\begin{defn}
Нека се извършват независими опити на Бернули с вероятност за успех $p$, където
\[
0<p<1.
\]
Нека $X$ е броят на неуспехите преди първия успех. Тогава $X$ има \emph{геометрично разпределение} с параметър $p$, означавано с
\[
X\sim\mathrm{Geo}(p),
\]
ако
\[
\mathbb{P}(X=k)=q^kp,
\qquad k=0,1,2,\ldots,
\]
където $q=1-p$.
\end{defn}

Събитието $\{X=k\}$ означава, че първите $k$ опита са неуспешни, а следващият опит е успешен. Поради независимостта вероятността му е
\[
q^kp.
\]

\begin{prop}
Вероятностите на геометричното разпределение са нормировани:
\[
\sum_{k=0}^{\infty}q^kp=1.
\]
\end{prop}

\begin{proof}
Тъй като $0<q<1$, редът е геометрична прогресия:
\[
\sum_{k=0}^{\infty}q^kp
=
p\sum_{k=0}^{\infty}q^k
=
p\frac{1}{1-q}.
\]
Понеже $1-q=p$, следва
\[
p\frac{1}{1-q}=\frac{p}{p}=1.
\]
\end{proof}

\begin{example}
Хвърляме правилен зар, докато се падне шестица. Ако $X$ е броят на неуспешните хвърляния преди първата шестица, то
\[
X\sim\mathrm{Geo}\left(\frac16\right).
\]
Например събитието $X=4$ означава, че са се получили четири нешестици, последвани от шестица, и има вероятност
\[
\mathbb{P}(X=4)
=
\left(\frac56\right)^4\frac16.
\]
\end{example}

\begin{supp}[Бележка]
В някои формулировки случайна величина се дефинира като броя на \emph{опитите} до първия успех. Тогава тя приема стойности $1,2,\ldots$ и се различава с единица от използваната тук величина. В настоящата глава следваме дефиницията, при която $X$ брои неуспехите преди първия успех; затова стойностите са $0,1,2,\ldots$ и
\[
\mathbb{P}(X=k)=q^kp.
\]
\end{supp}

\subsection{Пораждаща функция, математическо очакване и дисперсия}

За $X\sim\mathrm{Geo}(p)$ имаме
\[
\begin{aligned}
G_X(s)
&=\sum_{k=0}^{\infty}\mathbb{P}(X=k)s^k\\
&=\sum_{k=0}^{\infty}q^kps^k\\
&=p\sum_{k=0}^{\infty}(qs)^k.
\end{aligned}
\]
Като използваме сумата на геометрична прогресия, получаваме
\[
G_X(s)=\frac{p}{1-qs}.
\]

Първата производна е
\[
G_X'(s)=\frac{pq}{(1-qs)^2}.
\]
Следователно
\[
\mathbb{E}X
=
G_X'(1)
=
\frac{pq}{(1-q)^2}
=
\frac{pq}{p^2}
=
\frac{q}{p}.
\]

Втората производна е
\[
G_X''(s)=\frac{2pq^2}{(1-qs)^3}.
\]
Следователно
\[
G_X''(1)=\frac{2pq^2}{(1-q)^3}
=\frac{2q^2}{p^2}.
\]
Прилагаме формулата за дисперсията:
\[
\begin{aligned}
\mathbb{D}X
&=G_X''(1)+G_X'(1)-\bigl(G_X'(1)\bigr)^2\\
&=\frac{2q^2}{p^2}+\frac{q}{p}-\frac{q^2}{p^2}\\
&=\frac{q^2+pq}{p^2}\\
&=\frac{q(p+q)}{p^2}\\
&=\frac{q}{p^2}.
\end{aligned}
\]

\begin{keythm}
Ако
\[
X\sim\mathrm{Geo}(p),
\]
където $X$ е броят неуспехи преди първия успех, то
\[
\mathbb{E}X=\frac{q}{p},
\qquad
\mathbb{D}X=\frac{q}{p^2}.
\]
\end{keythm}

\begin{proof}
За $0<p\le1$ и $q=1-p$ пораждащата функция е $G(s)=p/(1-qs)$. Получаваме $G'(1)=pq/(1-q)^2=q/p$ и $G''(1)=2pq^2/(1-q)^3=2q^2/p^2$. По формулата за дисперсията
\[\mathbb D X=\frac{2q^2}{p^2}+\frac qp-\frac{q^2}{p^2}=\frac{q(q+p)}{p^2}=\frac q{p^2}.\]
При $p=1$ величината е тъждествено нула и формулите също са валидни.
\end{proof}

\section{Поасоново разпределение}

\subsection{Дефиниция и задачи, в които възниква}

Поасоновото разпределение се използва за броя на редки събития при голям брой възможности за настъпването им. То възниква като граничен случай на биномното разпределение, когато броят на опитите е голям, вероятността за успех във всеки опит е малка, а средният брой успехи остава краен.

\begin{defn}
Казваме, че случайната величина $X$ има \emph{Поасоново разпределение} с параметър $\lambda>0$, ако приема стойностите
\[
0,1,2,\ldots
\]
с вероятности
\[
\mathbb{P}(X=k)=\frac{e^{-\lambda}\lambda^k}{k!},
\qquad k=0,1,2,\ldots .
\]
Означаваме
\[
X\sim\mathrm{Po}(\lambda).
\]
\end{defn}

Параметърът $\lambda$ е средният брой настъпвания на разглежданото събитие.

\begin{example}
Дисплей има $1\,000\,000$ пиксела. Вероятността един пиксел да изгори е
\[
p=\frac{1}{1\,000\,000}.
\]
При $n=1\,000\,000$ възможности средният брой изгорели пиксели е
\[
\lambda=np=1.
\]
Броят $X$ на изгорелите пиксели може да се моделира с Поасоново разпределение с параметър $\lambda=1$. Например
\[
\mathbb{P}(X=2)
=
\frac{1^2e^{-1}}{2!}
=
\frac{1}{2e}.
\]
\end{example}

\begin{prop}
Вероятностите на Поасоновото разпределение са нормировани:
\[
\sum_{k=0}^{\infty}\frac{e^{-\lambda}\lambda^k}{k!}=1.
\]
\end{prop}

\begin{proof}
Използваме реда за експоненциалната функция:
\[
e^\lambda=\sum_{k=0}^{\infty}\frac{\lambda^k}{k!}.
\]
Тогава
\[
\sum_{k=0}^{\infty}\frac{e^{-\lambda}\lambda^k}{k!}
=
e^{-\lambda}
\sum_{k=0}^{\infty}\frac{\lambda^k}{k!}
=
e^{-\lambda}e^\lambda
=
1.
\]
\end{proof}

\subsection{Връзка с биномното разпределение}

Нека
\[
X\sim\mathrm{Bin}(n,p).
\]
Когато $n$ е голямо, $p$ е малко и
\[
np\longrightarrow\lambda,
\]
биномната вероятност
\[
\mathbb{P}(X=k)=\binom{n}{k}p^k(1-p)^{n-k}
\]
се стреми към
\[
\frac{e^{-\lambda}\lambda^k}{k!},
\qquad k=0,1,2,\ldots .
\]
Това обяснява защо Поасоновото разпределение е подходящ модел за редки успехи сред много на брой опити.

\subsection{Пораждаща функция, математическо очакване и дисперсия}

Нека
\[
X\sim\mathrm{Po}(\lambda).
\]
Тогава
\[
\begin{aligned}
G_X(s)
&=\sum_{k=0}^{\infty}\mathbb{P}(X=k)s^k\\
&=\sum_{k=0}^{\infty}
\frac{e^{-\lambda}\lambda^k}{k!}s^k\\
&=e^{-\lambda}\sum_{k=0}^{\infty}
\frac{(\lambda s)^k}{k!}.
\end{aligned}
\]
От реда за експоненциалната функция следва
\[
G_X(s)=e^{-\lambda}e^{\lambda s}
=e^{\lambda(s-1)}.
\]

Първата производна е
\[
G_X'(s)=\lambda e^{\lambda(s-1)}.
\]
Следователно
\[
\mathbb{E}X=G_X'(1)=\lambda.
\]

Втората производна е
\[
G_X''(s)=\lambda^2e^{\lambda(s-1)}.
\]
Следователно
\[
G_X''(1)=\lambda^2.
\]
Тогава
\[
\begin{aligned}
\mathbb{D}X
&=G_X''(1)+G_X'(1)-\bigl(G_X'(1)\bigr)^2\\
&=\lambda^2+\lambda-\lambda^2\\
&=\lambda.
\end{aligned}
\]

\begin{keythm}
Ако
\[
X\sim\mathrm{Po}(\lambda),
\]
то
\[
\mathbb{E}X=\lambda,
\qquad
\mathbb{D}X=\lambda.
\]
\end{keythm}

\begin{proof}
Пораждащата функция е $G(s)=e^{\lambda(s-1)}$. Следователно $G'(1)=\lambda$ и $G''(1)=\lambda^2$. От доказаните формули за моментите следва $\mathbb E X=\lambda$ и $\mathbb D X=\lambda^2+\lambda-\lambda^2=\lambda$.
\end{proof}

Равенството
\[
\mathbb{E}X=\mathbb{D}X=\lambda
\]
е характерна особеност на Поасоновото разпределение.

\section{Сравнение на разглежданите разпределения}

Трите разпределения описват различни въпроси, свързани с опити с успех и неуспех.

При биномното разпределение броят опити е фиксиран, а се броят успехите:
\[
X\sim\mathrm{Bin}(n,p),
\qquad
\mathbb{P}(X=k)=\binom{n}{k}p^kq^{n-k}.
\]
То е подходящо за въпроса: ``Колко успеха ще има в точно $n$ независими опита?''

При геометричното разпределение опитите продължават до първия успех, а се броят неуспехите преди него:
\[
X\sim\mathrm{Geo}(p),
\qquad
\mathbb{P}(X=k)=q^kp.
\]
То е подходящо за въпроса: ``Колко неуспеха ще има преди първия успех?''

При Поасоновото разпределение се брои броят на редки събития:
\[
X\sim\mathrm{Po}(\lambda),
\qquad
\mathbb{P}(X=k)=\frac{e^{-\lambda}\lambda^k}{k!}.
\]
То е подходящо, когато има много възможности за настъпване на събитие, всяка с малка вероятност, и средният брой настъпвания е $\lambda$.

Основните формули са:
\[
\begin{array}{c|c|c|c}
\text{Разпределение} & \text{Параметри} & \mathbb{E}X & \mathbb{D}X\\
\hline
\mathrm{Bin}(n,p) & n,\ p,\ q=1-p & np & npq\\
\mathrm{Geo}(p) & p,\ q=1-p & \dfrac{q}{p} & \dfrac{q}{p^2}\\
\mathrm{Po}(\lambda) & \lambda>0 & \lambda & \lambda
\end{array}
\]

Съответните пораждащи функции са:
\[
G_X(s)=
\begin{cases}
(q+ps)^n, & X\sim\mathrm{Bin}(n,p),\\[4pt]
\dfrac{p}{1-qs}, & X\sim\mathrm{Geo}(p),\\[10pt]
e^{\lambda(s-1)}, & X\sim\mathrm{Po}(\lambda).
\end{cases}
\]

\section{Изпитен фокус}

\begin{itemize}
\item Да се дефинира дискретна случайна величина и да се представи нейното разпределение чрез стойности $x_j$ и вероятности $p_j=\mathbb{P}(X=x_j)$.

\item Да се посочат условията
\[
p_j\geq0,
\qquad
\sum_jp_j=1,
\]
като израз на неотрицателността и нормираността на вероятностите.

\item Да се дефинира пораждащата функция
\[
G_X(s)=\sum_{k\geq0}\mathbb{P}(X=k)s^k
\]
и да се използват формулите
\[
G_X(1)=1,\qquad
\mathbb{E}X=G_X'(1),\qquad
\mathbb{D}X=G_X''(1)+G_X'(1)-\bigl(G_X'(1)\bigr)^2.
\]

\item Да се дефинира биномното разпределение като брой успехи в $n$ независими опита на Бернули:
\[
\mathbb{P}(X=k)=\binom{n}{k}p^kq^{n-k}.
\]
Да се изведе
\[
G_X(s)=(q+ps)^n,
\qquad
\mathbb{E}X=np,
\qquad
\mathbb{D}X=npq.
\]

\item Да се обясни геометричното разпределение като брой неуспехи преди първия успех:
\[
\mathbb{P}(X=k)=q^kp.
\]
Да се изведе
\[
G_X(s)=\frac{p}{1-qs},
\qquad
\mathbb{E}X=\frac{q}{p},
\qquad
\mathbb{D}X=\frac{q}{p^2}.
\]

\item Да се дефинира Поасоновото разпределение:
\[
\mathbb{P}(X=k)=\frac{e^{-\lambda}\lambda^k}{k!}.
\]
Да се изведе
\[
G_X(s)=e^{\lambda(s-1)},
\qquad
\mathbb{E}X=\lambda,
\qquad
\mathbb{D}X=\lambda.
\]

\item Да се разпознават типичните задачи: фиксиран брой опити за биномното разпределение, чакане до първи успех за геометричното разпределение и редки събития при много възможности за Поасоновото разпределение.

\item Да се възпроизведат идеите за проверка на нормираността: биномната формула на Нютон при биномното разпределение, сумата на геометрична прогресия при геометричното разпределение и редът за $e^x$ при Поасоновото разпределение.
\end{itemize}
